\section{Билет 33. Эффект Доплера. Формула частоты, когда приемник покоится, источник двигается вдоль одной прямой (с выводом). Обобщенная формула, когда источник и приемник двигаются вдоль одной прямой.}

\begin{definition}
\textbf{Эффектом Доплера} называется изменение частоты волн, регистрируемых приёмником, вследствие относительного движения источника волн и приёмника. Явление открыто австрийским физиком К. Доплером в 1842 г.
\par\noindent\textbf{Пример:} При приближении поезда к наблюдателю звук его гудка кажется более высоким, а при удалении --- более низким, чем у неподвижного источника.
\end{definition}

\begin{theorem}[Движущийся источник, неподвижный приёмник]
Пусть приёмник покоится, а источник звука движется вдоль прямой, соединяющей их, со скоростью $v_{\text{ист}}$. Скорость звука в среде --- $v$, частота источника --- $\nu_0$, период $T_0 = 1/\nu_0$.

\textbf{Случай 1: источник приближается к приёмнику.}
За время $T_0$ источник излучает одну волну. За это же время источник смещается на расстояние $v_{\text{ист}} T_0$ в сторону приёмника. Длина волны, регистрируемая неподвижным приёмником, уменьшается:
\begin{equation}
    \lambda = \frac{v T_0 - v_{\text{ист}} T_0}{1} = (v - v_{\text{ист}}) T_0 = \frac{v - v_{\text{ист}}}{\nu_0}.
\end{equation}
Частота, регистрируемая приёмником:
\begin{equation}
    \nu = \frac{v}{\lambda} = \nu_0 \frac{v}{v - v_{\text{ист}}}. \label{eq:doppler_source_approach}
\end{equation}
Так как знаменатель меньше $v$, частота возрастает: $\nu > \nu_0$.

\textbf{Случай 2: источник удаляется от приёмника.}
Аналогично, но источник «убегает» от своих волн:
\begin{equation}
    \lambda = (v + v_{\text{ист}}) T_0, \quad \nu = \nu_0 \frac{v}{v + v_{\text{ист}}}. \label{eq:doppler_source_recede}
\end{equation}
Частота уменьшается: $\nu < \nu_0$.
\end{theorem}

\begin{theorem}[Обобщённая формула эффекта Доплера]
Пусть и источник, и приёмник движутся вдоль одной прямой со скоростями $v_{\text{ист}}$ и $v_{\text{пр}}$ соответственно. Положительное направление оси $x$ --- от источника к приёмнику.

\textbf{Вывод:}
\begin{enumerate}
    \item Длина волны, излучаемая движущимся источником:
    \begin{equation}
        \lambda = \frac{v \mp v_{\text{ист}}}{\nu_0},
    \end{equation}
    где знак «минус» --- при сближении (источник «догоняет» свои волны), «плюс» --- при удалении.
    \item Частота, регистрируемая движущимся приёмником:
    \begin{equation}
        \nu = \frac{v \pm v_{\text{пр}}}{\lambda}.
    \end{equation}
    Знак «плюс» --- при сближении, «минус» --- при удалении.
\end{enumerate}
Подставляя $\lambda$, получаем \textbf{общую формулу эффекта Доплера}:
\begin{equation}
    \nu = \nu_0 \frac{v \pm v_{\text{пр}}}{v \mp v_{\text{ист}}}. \label{eq:doppler_general_33}
\end{equation}
\end{theorem}

\begin{property}[Правило знаков]
В формуле \eqref{eq:doppler_general_33}:
\begin{itemize}
    \item В числителе: знак «$+$» --- если приёмник \textbf{движется к источнику}, «$-$» --- если \textbf{удаляется}.
    \item В знаменателе: знак «$-$» --- если источник \textbf{движется к приёмнику}, «$+$» --- если \textbf{удаляется}.
\end{itemize}
Если скорости направлены не вдоль соединяющей прямой, в формулу подставляются проекции скоростей на эту прямую: $v_{\text{пр},x}$, $v_{\text{ист},x}$.
\end{property}

\begin{remark}
\textbf{Применение эффекта Доплера:}
\begin{enumerate}
    \item \textbf{Радары и лидары} --- измерение скорости автомобилей, самолётов, атмосферных потоков.
    \item \textbf{Медицинская диагностика} --- ультразвуковая допплерография кровотока.
    \item \textbf{Астрофизика} --- определение скорости звёзд и галактик по смещению спектральных линий (красное/синее смещение).
    \item \textbf{Навигация} --- спутниковые системы (GPS/ГЛОНАСС) учитывают доплеровский сдвиг для точного позиционирования.
\end{enumerate}
\end{remark}
